\subsection{双切片扭曲层析、完成化 Frobenius 塔与非阿贝尔奇性账本}

以下八条结论把单原子残差层析、cyclotomic 原子壳的指数透明性、prime-length/prime-power 完成化同余塔，以及 Lee--Yang 分支坐标与乘法账本的算术裂口压缩到同一条派生链上。除第一个层析命题采用自包含记号外，其余陈述均直接依托正文与附录中已审计的 local cumulant、completed Dwork congruence、Lee--Yang 三次数域与半圆维分离结论。

\begin{theorem}[单原子残差的双扭曲切片刚性层析]\label{thm:derived-atomic-two-slice-tomography}
\leanverified{paper\_derived\_atomic\_two\_slice\_tomography}
设 \(q_1,\dots,q_r\) 两两互素，记
\[
Q:=\prod_{i=1}^{r}q_i.
\]
设某一单原子差分通道在唯一指标 \(\ell\in [0,Q)\) 上非零，并对每个 \(i\) 记
\[
\ell_i\in [0,q_i),\qquad \ell_i\equiv \ell \pmod{q_i}.
\]
再固定某个 \(i_0\in\{1,\dots,r\}\) 及两条不同扭曲切片 \(u_1,u_2>0\)，并设在该切片上有
\[
c_{i_0}^{(j)}
:=
[z^{\ell_{i_0}}]\Delta a_{i_0,q_{i_0}}(z;u_j)
=
m\,u_j^{E}
\qquad (j=1,2).
\]
则三元组 \((\ell,E,m)\) 被有限数据
\[
\bigl((\ell_i)_{i=1}^{r},\,c_{i_0}^{(1)},\,c_{i_0}^{(2)}\bigr)
\]
唯一决定，且有闭式恢复公式
\[
E=\frac{\log\!\bigl(c_{i_0}^{(2)}/c_{i_0}^{(1)}\bigr)}{\log(u_2/u_1)},
\qquad
m=c_{i_0}^{(1)}u_1^{-E}.
\]
\end{theorem}

\begin{proof}
由中国剩余定理，模 \(q_i\) 的 residue 数据 \((\ell_i)_{i=1}^{r}\) 唯一确定区间 \([0,Q)\) 中的整数 \(\ell\)。在 \(\ell\) 已知后，两个切片系数满足
\[
\frac{c_{i_0}^{(2)}}{c_{i_0}^{(1)}}
=
\left(\frac{u_2}{u_1}\right)^E.
\]
因 \(u_1\neq u_2\)，故上式唯一给出
\[
E=\frac{\log(c_{i_0}^{(2)}/c_{i_0}^{(1)})}{\log(u_2/u_1)}.
\]
再由
\(
c_{i_0}^{(1)}=m\,u_1^E
\)
反解得
\(
m=c_{i_0}^{(1)}u_1^{-E}
\)。
证毕。
\end{proof}

\begin{corollary}[单切片不可识别，而双切片已达最小完备]\label{cor:derived-atomic-one-vs-two-slice-identifiability}
\leanverified{paper\_derived\_atomic\_one\_vs\_two\_slice\_identifiability}
在定理 \ref{thm:derived-atomic-two-slice-tomography} 的设定下，若 \(\ell\) 已由 residue 支撑数据唯一恢复，而只给定单一切片 \(u>0\) 上的一个系数
\[
c=[z^{\ell_{i_0}}]\Delta a_{i_0,q_{i_0}}(z;u)=m\,u^E,
\]
则在不额外施加离散先验、上界先验或外部归一化条件时，\((m,E)\) 不可唯一识别；而任取两条不同切片 \(u_1\neq u_2\)，便已足以唯一恢复全部连续原子参数。
\end{corollary}

\begin{proof}
若仅知 \(c=m\,u^E\) 且 \(u\neq 1\)，则对任意 \(t\in\RR\) 置
\[
m_t:=m\,u^{-t},
\qquad
E_t:=E+t,
\]
便有
\[
m_tu^{E_t}=m\,u^E=c.
\]
故单切片只施加一条乘法约束，而不能唯一标定 \((m,E)\)。若 \(u=1\)，则 \(c=m\)，此时 \(E\) 完全不可见。两切片的唯一恢复已由定理 \ref{thm:derived-atomic-two-slice-tomography} 给出，故“双切片充足、单切片不足”成立。证毕。
\end{proof}

\begin{theorem}[有限 cyclotomic 手术对正指数 residue 不变量完全透明]\label{thm:derived-finite-cyclotomic-surgery-transparent}
设某一 core residue 生成函数 \(G^{\mathrm{core}}(z)\) 经有限次单原子增补后变为
\[
G^{(N)}(z)
=
G^{\mathrm{core}}(z)
+
\sum_{\nu=1}^{N}\frac{2z^{\ell_\nu}}{1-z^{M_\nu}},
\qquad M_\nu\ge 2.
\]
记 \(\rho^{\mathrm{core}}\) 为 \(G^{\mathrm{core}}(z)\) 所决定的指数半径，\(\rho^{(N)}\) 为增补后的对应半径；并对某个固定 Perron 标尺 \(\lambda>1\) 设
\[
\theta^{\mathrm{core}}
:=
\frac{\log \rho^{\mathrm{core}}}{\log\lambda},
\qquad
\theta^{(N)}
:=
\frac{\log \rho^{(N)}}{\log\lambda}.
\]
则恒有
\[
\rho^{(N)}=\max\{1,\rho^{\mathrm{core}}\},
\qquad
\theta^{(N)}=\max\{0,\theta^{\mathrm{core}}\}.
\]
特别地，对任意阈值 \(\tau>0\)，
\[
\theta^{(N)}>\tau
\iff
\theta^{\mathrm{core}}>\tau.
\]
因此，任何有限个 cyclotomic 原子壳都不能修复、制造或改变任何真正正指数的 twisted 增长；尤其不可能修复平方根级门槛的失效。
\end{theorem}

\begin{proof}
每个附加项
\[
\frac{2z^{\ell_\nu}}{1-z^{M_\nu}}
\]
的极点全部位于 \(M_\nu\)-次单位根集合上，故其全部极点模长都等于 \(1\)。因此 \(G^{(N)}(z)\) 的一切模长不等于 \(1\) 的奇点与 \(G^{\mathrm{core}}(z)\) 完全相同；新添的奇点只能出现在单位圆上。于是，决定指数增长的最小模奇点若位于单位圆内，则只能来自 core；若 core 不提供模长小于 \(1\) 的奇点，则新模型的最小模奇点只能落在 \(|z|=1\)。故
\[
\rho^{(N)}=\max\{1,\rho^{\mathrm{core}}\}.
\]
再取 \(\log/\log\lambda\) 即得
\[
\theta^{(N)}=\max\{0,\theta^{\mathrm{core}}\}.
\]
阈值不变性随即成立。单原子情形正是推论 \ref{cor:xi-time-part50dc-full-periodic-residue-core-cyclotomic-splitting} 的直接特例。证毕。
\end{proof}

\begin{theorem}[prime-length 完成化轨道数据的 Frobenius 塌缩]\label{thm:derived-prime-length-completed-frobenius-collapse}
\leanverified{paper\_derived\_prime\_length\_completed\_frobenius\_collapse}
沿用 completed ghost 多项式 \(\widehat a_n(s)\) 与 Chebyshev--Adams 多项式 \(C_p(s)\) 的记号。对任意素数 \(p\)，有双重同余
\[
\widehat a_p(s)\equiv C_p(s)\equiv s^p \pmod p.
\]
从而，长度 \(p\) 的完成化轨道数据在模 \(p\) 上并不携带任何超出一次坐标 \(s\) 的新信息；它只是 Frobenius 自同态
\[
s\longmapsto s^p
\]
的重写。
\end{theorem}

\begin{proof}
推论 \ref{cor:completed-teichmueller-congruence} 已给出
\[
\widehat a_p(s)\equiv s^p \pmod p.
\]
而其证明正是先由
\(
\widehat a_p(s)\equiv C_p(s)\pmod p
\)
再结合
\(
C_p(s)\equiv s^p\pmod p
\)
得到。于是
\[
\widehat a_p(s)\equiv C_p(s)\equiv s^p\pmod p.
\]
因此任意 \(\FF_p\)-多项式可观测量 \(F\) 都满足
\[
F(\widehat a_p(s))=F(s^p)=F(s)^p,
\]
即 length-\(p\) 的 completed 数据完全因子化为 Frobenius 像。证毕。
\end{proof}

\begin{theorem}[prime-power 长度的 Dwork--Frobenius 塔]\label{thm:derived-prime-power-dwork-frobenius-tower}
\leanverified{paper\_derived\_prime\_power\_dwork\_frobenius\_tower}
沿用同一完成化记号。对任意素数 \(p\) 与整数 \(k\ge 1\)，有精确整系数同余
\[
\widehat a_{p^k}(s)-\widehat a_{p^{k-1}}\!\bigl(C_p(s)\bigr)\in p^k\ZZ[s].
\]
因此
\[
\widehat a_{p^k}(s)\equiv \widehat a_{p^{k-1}}\!\bigl(C_p(s)\bigr)\pmod{p^k}.
\]
进一步，模 \(p\) 约化后整条 prime-power 塔严格塌缩为 Frobenius 迭代：
\[
\widehat a_{p^k}(s)\equiv \widehat a_{p^{k-1}}(s^p)\equiv s^{p^k}\pmod p.
\]
\end{theorem}

\begin{proof}
第一条正是定理 \ref{thm:discussion-completed-dwork-chebyshev} 的整系数重述。对模 \(p\) 约化，由推论 \ref{cor:completed-dwork-frobenius-tower} 立得
\[
\widehat a_{p^k}(s)\equiv \widehat a_{p^{k-1}}(s^p)\pmod p.
\]
再对 \(k\) 归纳，并以定理 \ref{thm:derived-prime-length-completed-frobenius-collapse} 的
\(
\widehat a_p(s)\equiv s^p\pmod p
\)
为基步，即得
\[
\widehat a_{p^k}(s)\equiv s^{p^k}\pmod p.
\]
证毕。
\end{proof}

\begin{theorem}[正则 germ 与奇性分支的算术相分裂]\label{thm:derived-regular-singular-arithmetic-splitting}
沿用命题 \ref{prop:real-input-40-output-cumulants-closed} 的输出位势记号 \(P(\theta)\)，并沿用定理 \ref{thm:killo-leyang-kappa-square-cubic-field} 的 Lee--Yang 分支坐标
\[
u:=\kappa_\star^2.
\]
定义
\[
K_{\mathrm{reg}}:=\QQ(\sqrt5),
\qquad
K_{\mathrm{sing}}:=\QQ(u).
\]
则：
\begin{enumerate}
  \item 对一切 \(k\ge 1\)，局部累积量都满足
  \[
  P^{(k)}(0)\in K_{\mathrm{reg}};
  \]
  \item \(u\) 满足不可约三次方程
  \[
  324u^3-540u^2+333u-74=0,
  \]
  其判别式为
  \[
  -2^6\cdot 3^9\cdot 37,
  \]
  因而 \(K_{\mathrm{sing}}/\QQ\) 为非 Galois 三次域，其分裂域 Galois 群为 \(S_3\)；
  \item 二者交域严格平凡：
  \[
  K_{\mathrm{reg}}\cap K_{\mathrm{sing}}=\QQ.
  \]
\end{enumerate}
换言之，正则扰动 germ 与奇性分支数据在算术上严格横截：前者完全落在二次数域，后者则逃逸到一个非阿贝尔控制的三次扩张；不存在任何二次域参数化可同时吸收这两类信息。
\end{theorem}

\begin{proof}
命题 \ref{prop:real-input-40-output-cumulants-closed} 给出全部局部累积量
\[
P^{(k)}(0)\in \QQ(\sqrt5)=K_{\mathrm{reg}}
\qquad (k\ge 1).
\]
另一方面，定理 \ref{thm:killo-leyang-kappa-square-cubic-field} 已证明 \(u\) 满足所示不可约三次方程，且其判别式非平方；故 \(K_{\mathrm{sing}}=\QQ(u)\) 为非 Galois 三次域，而其分裂域 Galois 群为 \(S_3\)。

最后，交域次数同时整除
\[
[K_{\mathrm{reg}}:\QQ]=2,
\qquad
[K_{\mathrm{sing}}:\QQ]=3,
\]
故只能为 \(1\)。于是
\[
K_{\mathrm{reg}}\cap K_{\mathrm{sing}}=\QQ.
\]
若存在二次域统一化可同时吸收两类信息，则必有 \(u\in K_{\mathrm{reg}}\)，与 \([K_{\mathrm{sing}}:\QQ]=3\) 矛盾。证毕。
\end{proof}

\begin{proposition}[Lee--Yang 分支坐标不能被任何有限阿贝尔 residue 账本精确承载]\label{prop:derived-leyang-no-finite-abelian-ledger}
沿用定理 \ref{thm:derived-regular-singular-arithmetic-splitting} 的记号。不存在有限阿贝尔扩张 \(L/\QQ\) 使
\[
u=\kappa_\star^2\in L.
\]
因此，一切纯 cyclotomic、Teichm\"uller 或有限阿贝尔 residue 账本，都不可能对 Lee--Yang 分支坐标给出精确封闭表达。
\end{proposition}

\begin{proof}
若存在有限阿贝尔扩张 \(L/\QQ\) 使 \(u\in L\)，则 \(\QQ(u)\) 是阿贝尔 Galois 扩张 \(L/\QQ\) 的中间域。由于阿贝尔群的任意子群都是正规子群，故一切中间域都必须为 \(\QQ\) 上的 Galois 扩张。于是 \(\QQ(u)/\QQ\) 必为 Galois。

但定理 \ref{thm:derived-regular-singular-arithmetic-splitting} 已表明 \(\QQ(u)\) 是非 Galois 三次域，矛盾。故不存在这样的有限阿贝尔扩张 \(L\)。证毕。
\end{proof}

\begin{theorem}[乘法对象的实现维与同态维出现最大裂口]\label{thm:derived-multiplicative-implementation-hom-gap}
\leanverified{paper\_derived\_multiplicative\_implementation\_hom\_gap}
记命题 \ref{prop:killo-implementation-vs-structural-half-circle-dimension} 中定义的两类半圆维为
\[
\cdim^{\mathrm{impl}}(X),
\qquad
\cdim^{\mathrm{hom}}_{+}(M).
\]
则对乘法幺半群 \(\NN_{\times}\) 有精确等式
\[
\cdim^{\mathrm{impl}}(\NN_{\times})=\frac12,
\qquad
\cdim^{\mathrm{hom}}_{+}(\NN_{\times})=+\infty.
\]
亦即：作为集合，\(\NN_{\times}\) 可以压缩到单寄存器；作为乘法结构，它却不能被任何有限维加法寄存器忠实同态化。实现压缩与结构压缩在此达到最强分离。
\end{theorem}

\begin{proof}
这正是命题 \ref{prop:killo-implementation-vs-structural-half-circle-dimension} 的后两条精确等式。第一式来自 \(\NN_{\times}\) 作为可数无限集可单射嵌入 \(\NN\)；第二式则由推论 \ref{cor:killo-no-finite-additive-register-linearization} 强制：不存在
\[
\NN_{\times}\hookrightarrow (\NN^k,+)
\]
的单射幺半群同态。故任何声称以有限加法寄存器完整保留乘法结构的账本，都必在“忠实”“同态”“有限维”三者中至少失去一项。证毕。
\end{proof}

综上，单原子残差参数在两条扭曲切片上已被刚性锁定；有限个 cyclotomic 补丁对一切真正正指数现象完全透明；prime-length 与 prime-power 完成化数据在模 \(p\) 上塌缩为 Frobenius 塔；正则扰动 germ 仅活在 \(\QQ(\sqrt5)\)，而奇性分支坐标却逃入一个 \(S_3\)-控制的非阿贝尔三次域；乘法结构的有限实现压缩与有限同态压缩之间则存在不可弥合的量级裂口。因而，可见层连续读出、有限 residue 补丁、模 \(p\) 同余护栏与乘法账本压缩并不共享同一算术层级；真正的奇性信息只在严格高于 Frobenius--cyclotomic 世界的非阿贝尔层中才开始出现。
